\documentclass[10pt]{beamer}
\usepackage{hyperref} % Create hyperlinks in your document
\usepackage[style=authortitle]{biblatex} % For citation
\renewcommand*{\bibfont}{\tiny}
\addbibresource{bibliography.bib}
%\usepackage{newtxmath} % If you don't have stix2 font, uncomment this line and comment the next line
%\usepackage{stix2} % Use stix2 font for math
\usepackage[no-math]{fontspec}
\usepackage{xeCJK} % Chinese font setting
\usepackage{tcolorbox} % Coloured boxes, for LATEX examples and theorems, etc
\tcbuselibrary{skins, breakable, theorems}
\usepackage{url} % Inserting urls
\usepackage{caption} % Customising captions in floating environments
\usepackage{verbatim}
\usepackage{multicol} % Multi-columns
\usepackage{booktabs} % Better tables
\usepackage{tikz} % For plot
\usetikzlibrary{angles, arrows.meta, quotes, calc, intersections}
\usepackage{pgfplots}
\pgfplotsset{compat=1.16}
\usepackage{soul} % For smallcaps/striking out/highlighting
\sethlcolor{cuhklightyellow}
\usepackage{minted} % For source code
\usepackage{ifplatform} % For identifying OS platforms
\usepackage[vlined, ruled, linesnumbered]{algorithm2e} % For algorithms
\usepackage{algorithmic}
\usepackage{amsmath} % For math symbols
\usepackage{amssymb}
\usepackage{stmaryrd}
\usepackage{pifont}
\usepackage{varwidth}
\usepackage{outlines} % For simplifying nested itemized/enumerate lists
\usepackage{braket} % For quantum Dirac notation
\usepackage{derivative} % For derivatives

% Set minted style and font size, used for displaying source code
\setminted[rust]{frame=lines, breaklines, xleftmargin=16pt, linenos, fontsize=\tiny}
\setminted[c]{frame=lines, breaklines, xleftmargin=16pt, linenos, fontsize=\tiny}
\setminted[python]{frame=lines, breaklines, xleftmargin=16pt, linenos, fontsize=\tiny}
\setminted[bash]{frame=lines, breaklines, xleftmargin=16pt, linenos, fontsize=\tiny}
% For inline code, use the current font size
\makeatletter
\newcommand{\currentfontsize}{\fontsize{\f@size}{\f@baselineskip}\selectfont}
\makeatother
\setmintedinline{fontsize=\currentfontsize}

\definecolor{cuhkpurple}{RGB}{117,15,109}
\definecolor{cuhkyellow}{RGB}{221,163,0}
\definecolor{cuhkblackyellow}{RGB}{153,102,0}
\definecolor{cuhklightyellow}{RGB}{244,223,176}

\newtcbtheorem{theo}{Theorem}%
  {enhanced, breakable,
    colback = white, colframe = cuhkpurple, colbacktitle = cuhkpurple,
    attach boxed title to top left = {yshift = -2mm, xshift = 5mm},
    boxed title style = {sharp corners},
    fonttitle = \sffamily\bfseries}{th}

\iflinux
  \setCJKmainfont{Source Han Sans CN}
  \setCJKsansfont{Source Han Sans CN}
\fi
\ifmacosx
  \setCJKmainfont{Source Han Sans}
  \setCJKsansfont{Source Han Sans}
\fi

\usefonttheme{professionalfonts}
\setmainfont[Mapping=tex-text]{Helvetica Neue}
\setsansfont{Helvetica Neue}
\setmonofont{Source Code Pro}

\useoutertheme{infolines}
\setbeamertemplate{navigation symbols}{}
\usetheme[height=10mm]{Rochester}
\usecolortheme[RGB={117,15,109}]{structure}
\setbeamercolor{block title example}{use=structure,fg=white, bg=cuhkyellow}
\setbeamercolor{block body example}{use=structure,fg=black, bg=cuhklightyellow}
\setbeamercolor{alerted text}{fg=cuhkpurple, bg=cuhklightyellow}
\setbeamercolor{frametitle}{ fg=white, bg=cuhkpurple }

\newcommand{\boxalert}[1]{{
    \usebeamercolor{alerted text}\colorbox{bg}{\alert{#1}}
}}

\newcommand{\cmark}{\ding{51}}
\newcommand{\xmark}{\ding{55}}

% Define some useful abbreviations
\def\ZZ{{\mathbb Z}}
\def\RR{{\mathbb R}}
\def\CC{{\mathbb C}}
\def\QQ{{\mathbb Q}}
\def\EE{{\mathbb E}}
\newcommand{\tr}{{\rm tr}}


%% Logo
\titlegraphic{
    \includegraphics[width=4cm]{The-Chinese-University-of-Hong-Hong-logo.eps}
}

